
\sgaStableTarget{sga6.e10.sec.1.3}\subsection*{1.3.}
Recall (IV 2.10) that the bifunctor
@@M0@@
where @@M1@@ is a locally free @@M2@@-module and @@M3@@ a coherent @@M4@@-module, induces a composition law
\sgaStableTarget{sga6.e10.eq.1-3-1}@@M5@@
that makes @@M6@@ a unital @@M7@@-module. The key result of the present exposé is the following

\sgaStableTarget{sga6.e10.thm.1-3-2}\textbf{Theorem 1.3.2.} \uline{Let} @@M8@@ \uline{be a noetherian scheme. For all} @@M9@@, \uline{the composition law} (1.3.1) \uline{induces by restriction a composition law}
\sgaStableTarget{sga6.e10.eq.1-3-3}@@M10@@

\sgaStableTarget{sga6.e10.cor.1-3-3}\textbf{Corollary 1.3.3.} \uline{If} @@M11@@ \uline{is a noetherian scheme of dimension} @@M12@@\uline{, then the canonical homomorphism}
@@M13@@
\uline{defined by} @@M14@@ (@@M15@@\uline{) induces for every} @@M16@@ \uline{a homomorphism}
@@M17@@

\emph{Proof:} First let us show that (1.3.3), for @@M18@@ and for \emph{every} integral closed subscheme of dimension @@M19@@, implies (1.3.2) for @@M20@@. By (1.1.4), it is enough to show that for every @@M21@@ and for \emph{every} integral closed subscheme @@M22@@ of dimension @@M23@@ one has

